\section{Representation, correction and combination}\label{sec:results-representation}

\paragraph{Concatenated features.}
Table~\ref{tab:v2-replication} gives the three-member width-2000 experiment of
Section~\ref{sec:scaling} split by split, with the networks trained on one GPU.

\begin{table}[p]
\centering
% Generated by code/format_publication_tables.py; numerical cells preserved.
{\fontsize{9}{11}\selectfont
\setlength{\tabcolsep}{2.8pt}
\begin{tabular}{llrrrr}
\toprule
seed & head & radiance [\%] & median [pp] & p95 [pp] & gain [pp] \\
\midrule
104 & feature kernel & 0.0924 & 0.044 & 9.25 & -- \\
104 & concatenated features & 0.0810 & 0.025 & 6.25 & +0.0114 \\
104 & campaign stack & 0.0781 & 0.024 & 12.77 & -- \\
104 & network + residual & 0.1843 & 0.115 & 54.71 & -- \\
105 & feature kernel & 0.0611 & 0.031 & 7.31 & -- \\
105 & concatenated features & 0.0508 & 0.028 & 6.02 & +0.0103 \\
105 & campaign stack & 0.0475 & 0.024 & 14.50 & -- \\
105 & network + residual & 0.1349 & 0.087 & 42.61 & -- \\
106 & feature kernel & 0.0789 & 0.043 & 7.67 & -- \\
106 & concatenated features & 0.0702 & 0.035 & 6.28 & +0.0087 \\
106 & campaign stack & 0.0687 & 0.034 & 16.29 & -- \\
106 & network + residual & 0.1608 & 0.108 & 51.84 & -- \\
107 & feature kernel & 0.0731 & 0.042 & 5.82 & -- \\
107 & concatenated features & 0.0732 & 0.044 & 5.41 & -0.0002 \\
107 & campaign stack & 0.0833 & 0.043 & 13.84 & -- \\
107 & network + residual & 0.1476 & 0.103 & 46.58 & -- \\
mean & feature kernel & 0.0764$\pm$0.0130 & 0.040$\pm$0.006 & 7.51$\pm$1.41 & -- \\
mean & concatenated features & 0.0688$\pm$0.0128 & 0.033$\pm$0.009 & 5.99$\pm$0.40 & 0.0076$\pm$0.0053 \\
mean & campaign stack & 0.0694$\pm$0.0158 & 0.031$\pm$0.009 & 14.35$\pm$1.48 & -- \\
mean & network + residual & 0.1569$\pm$0.0211 & 0.103$\pm$0.012 & 48.93$\pm$5.40 & -- \\
\bottomrule
\end{tabular}}

\caption{Kernel on concatenated features, three width-2000 networks, splits 104--107. Radiance
is a mean relative spectral error in percent; reflectance median and 95th-percentile errors are
in percentage points. The last column is the paired radiance gain of concatenated features over
the features of one network. The mean row averages the four within-split quantities.}
\label{tab:v2-replication}
\end{table}

The concatenated-feature kernel has mean radiance error $0.0688\%$, against $0.0764\%$ for the
features of one network; the paired gain is $0.0076\pm0.0053$ percentage points, positive on
three splits and $-0.0002$ on the fourth. Its mean within-split 95th-percentile reflectance error
is $5.99$ points, against $7.51$ for single-member features and $14.35$ for the stack. With the
ten width-512 splits of Section~\ref{sec:scaling} this makes two settings, with different
networks on overlapping splits, in which concatenation lowers both the radiance error and the
reflectance tail.

\paragraph{Kernel choice on the state and on learned features.}
Table~\ref{tab:v2-multikernel} covers four components on seeds 101--103. The spectral mixture
and the Mat\'ern-5/2 kernel have similar mean errors on $Y_1$ ($0.089\%$), and the spectral
mixture has the lowest mean errors of the raw-state kernels on $Y_2$ and $Y_3$ ($0.237\%$ and
$0.099\%$). The first-order additive kernel is much less accurate in this specification.
Albedo errors for the nonadditive kernels lie between about $2.3\%$ and $3.5\%$.

\begin{table}[p]
\centering
% Generated by code/format_publication_tables.py; numerical cells preserved.
{\fontsize{9}{11}\selectfont
\setlength{\tabcolsep}{3pt}
\begin{tabular}{lrrrr}
\toprule
head & $Y_1$ [\%] & $Y_2$ [\%] & $Y_3$ [\%] & $Y_4$ [\%] \\
\midrule
Mat\'ern-1/2 & 0.179$\pm$0.008 & 0.672$\pm$0.003 & 0.418$\pm$0.012 & 2.326$\pm$0.103 \\
Mat\'ern-3/2 & 0.092$\pm$0.002 & 0.391$\pm$0.008 & 0.207$\pm$0.016 & 2.481$\pm$0.043 \\
Mat\'ern-5/2 & 0.089$\pm$0.017 & 0.434$\pm$0.007 & 0.203$\pm$0.003 & 2.746$\pm$0.082 \\
Gaussian & 0.092$\pm$0.002 & 0.525$\pm$0.113 & 0.228$\pm$0.051 & 3.502$\pm$0.501 \\
spectral mixture (4) & 0.089$\pm$0.010 & 0.237$\pm$0.014 & 0.099$\pm$0.013 & 3.468$\pm$0.263 \\
NNGP, ReLU & 0.175$\pm$0.001 & 0.674$\pm$0.015 & 0.350$\pm$0.009 & 2.548$\pm$0.057 \\
NTK, ReLU & 0.263$\pm$0.002 & 1.055$\pm$0.015 & 0.621$\pm$0.014 & 2.321$\pm$0.048 \\
additive Mat\'ern-5/2 & 4.091$\pm$0.093 & 20.266$\pm$1.035 & 11.430$\pm$0.237 & 4.413$\pm$0.077 \\
ARD by kernel flow & 0.322$\pm$0.015 & 0.257$\pm$0.019 & 0.199$\pm$0.005 & 2.418$\pm$0.114 \\
ARD by marginal likelihood & 0.281$\pm$0.039 & 0.341$\pm$0.042 & 0.183$\pm$0.011 & 2.702$\pm$0.027 \\
convex kernel sum & 0.087$\pm$0.007 & 0.253$\pm$0.020 & 0.099$\pm$0.013 & 2.557$\pm$0.099 \\
network & 0.819$\pm$0.027 & 2.289$\pm$0.019 & 0.910$\pm$0.057 & 3.482$\pm$0.085 \\
network + residual kernel & 0.523$\pm$0.007 & 1.337$\pm$0.026 & 0.588$\pm$0.035 & 3.055$\pm$0.121 \\
features, Mat\'ern-5/2 & 0.105$\pm$0.006 & 0.211$\pm$0.016 & 0.114$\pm$0.008 & 2.303$\pm$0.216 \\
features, Gaussian & 0.126$\pm$0.007 & 0.262$\pm$0.018 & 0.138$\pm$0.014 & 2.426$\pm$0.214 \\
features, NTK & 0.152$\pm$0.005 & 0.307$\pm$0.015 & 0.156$\pm$0.007 & 2.150$\pm$0.059 \\
features, ARD & 0.116$\pm$0.007 & 0.214$\pm$0.014 & 0.111$\pm$0.004 & 2.217$\pm$0.207 \\
coordinate selection & 0.090$\pm$0.010 & 0.360$\pm$0.008 & 0.105$\pm$0.010 & 2.797$\pm$0.447 \\
convex stack of the heads & 0.066$\pm$0.003 & 0.156$\pm$0.010 & 0.077$\pm$0.004 & 1.956$\pm$0.037 \\
ridge stack of the heads & 0.087$\pm$0.005 & 0.288$\pm$0.031 & 0.087$\pm$0.005 & 2.310$\pm$0.071 \\
\bottomrule
\end{tabular}}

\caption{Ten input-space kernels, feature kernels and combination rules on the EMIT state
inputs, with rank-64 outputs. Entries are component mean relative spectral errors in percent,
averaged over seeds 101--103 with sample standard deviations. The network has width 512 and depth
3 with a 150-epoch budget. Kernels include the Mat\'ern family, Gaussian, spectral mixture
\cite{wilsonadams}, NNGP \cite{chosaul,leenngp}, NTK \cite{ntk}, additive kernel
\cite{additivegp}, and learned input scales \cite{owhadiyoo,rw}.}
\label{tab:v2-multikernel}
\end{table}

The Mat\'ern-5/2 feature kernel gives $0.105\%$, $0.211\%$, $0.114\%$ and $2.303\%$ on
$Y_1$--$Y_4$, better than the best raw-state kernel on some components and not on others. The
convex stack has the lowest component means in this comparison ($0.066\%$, $0.156\%$, $0.077\%$,
$1.956\%$), which is consistent with its large reflectance tails elsewhere: the component metric
does not weight errors the way the inverse does.

\paragraph{Network regularization.}
Table~\ref{tab:v2-regsuite} compares a control, a weight-decay grid and an early-stopping grid on
three seeds. The validation-selected weight decay matches the control's component mean
($2.27\%$), and the selected early stopping is less accurate and more variable
($2.72\pm0.45\%$). Residual correction brings the component means to $1.88$--$1.95\%$ and
feature kernels to $0.95$--$0.99\%$, with radiance errors near $0.15\%$, whichever
regularization the network used. The reference row is the isotropic kernel fitted in the same runs.
Under this budget, regularizing the network changes little; the kernel stage does the work.

\begin{table}[t]
\centering
% Generated by code/format_publication_tables.py; numerical cells preserved.
{\fontsize{9}{11}\selectfont
\setlength{\tabcolsep}{3pt}
\begin{tabular}{lrrr}
\toprule
head & \thead{alone\\component / radiance} & \thead{residual correction\\component / radiance} & \thead{feature kernel\\component / radiance} \\
\midrule
isotropic kernel & \shortstack{0.941$\pm$0.109\\0.284$\pm$0.078} & -- & -- \\
network, control & \shortstack{2.266$\pm$0.031\\0.809$\pm$0.029} & \shortstack{1.903$\pm$0.055\\0.617$\pm$0.033} & \shortstack{0.955$\pm$0.038\\0.153$\pm$0.052} \\
network, weight decay & \shortstack{2.269$\pm$0.027\\0.820$\pm$0.014} & \shortstack{1.880$\pm$0.072\\0.618$\pm$0.040} & \shortstack{0.948$\pm$0.048\\0.154$\pm$0.048} \\
network, early stopping & \shortstack{2.716$\pm$0.445\\1.293$\pm$0.373} & \shortstack{1.949$\pm$0.073\\0.789$\pm$0.127} & \shortstack{0.986$\pm$0.032\\0.147$\pm$0.054} \\
\bottomrule
\end{tabular}}

\caption{Regularization comparison, seeds 101--103. Each cell gives the mean component
relative error [percent] / radiance relative error [percent], with sample standard deviations.
The width-512, depth-3 network has a 150-epoch cap, and validation selects among three runs per
arm. The weight-decay grid is $10^{-5},10^{-4},10^{-3}$ and the early-stopping patience grid 3, 5
and 10 epochs \cite{adamw,prechelt}. The reference row is the isotropic kernel fitted in the same
runs.}
\label{tab:v2-regsuite}
\end{table}

\paragraph{Perturbation and ridge-path diagnostics.}
The tables of Section~\ref{app:diagnostics} give the perturbation diagnostics for the state
inputs and the learned features. The residual-action bound of Proposition~\ref{prop:perturb} is
often much smaller than the spectrum-only bound, yet it can still exceed the actual movement by a
factor of tens, thousands or more: positive definiteness makes the bound applicable without
making it a useful tolerance.

Table~\ref{tab:v2-ridge} reports a 4{,}000-row fit of four components at seed 101. At the
smallest ridge on input-space $Y_1$ the effective degrees of freedom are 3{,}841 and the noise
factor $0.924$. The same-half and split-half principal-component errors are $0.079\%$ and
$2.465\%$, and the physical-unit test errors of the corrected and standalone kernels $0.2168\%$
and $0.1754\%$; these columns use different metrics and samples and are not to be subtracted. At
$\gamma=10^{-2}$ the corresponding test errors are $0.5095\%$ and $5.7578\%$: the fitted mean
keeps its predictive information when the kernel correction is strongly regularized.

\section{A synthetic check of the transfer theorem}\label{sec:results-synthetic}

Table~\ref{tab:v2-synthetic} checks the construction of Proposition~\ref{prop:attained} on
manufactured data, two million samples for each $\beta\in\{0.5,1,2,3,4\}$. The average-error
slopes are $0.200$, $0.333$, $0.498$, $0.597$ and $0.664$, against the limits $0.200$, $0.333$,
$0.500$, $0.600$ and $0.667$. For the uniform family the finite-ladder slope is compared with the
analytic curve on the same ladder: at $\beta=2$ the measured slope is $1.721$ and the analytic one
$1.724$, while the limit is two. The largest relative Monte Carlo discrepancy is about $1.01\%$ for
the uniform family at $\beta=2$ and below $0.053\%$ for the profile family against its quadrature
reference, and no pointwise violation occurs among the ten million sampled transmissions.

\section{The matched benchmark and the transfer contrast}\label{sec:results-benchmarks}

\paragraph{Correction-coefficient benchmark.}
Table~\ref{tab:v2-pkanrtm} scores our families under the standard and out-of-distribution
protocols of \cite{pkan}, with and without the low-fidelity 6S coefficients as inputs, against the
published values of the reference model, which we did not retrain.

The state-only residual-corrected network reaches standard-split RMSE $0.00581\pm0.00004$, MAE
$0.00151$, uniform-average $R^2=0.9952$ and SMAPE $3.54\%$, against published values $0.00619$,
$0.00186$, $0.9947$ and $3.91\%$; with 6S inputs its means are $0.00578$, $0.00148$, $0.9954$ and
$3.19\%$. On the out-of-distribution split the state-only means are $0.00620$, $0.00252$, $0.9967$
and $3.37\%$, against published $0.01022$, $0.00544$, $0.9942$ and $5.40\%$, and with 6S inputs
$0.00596$, $0.00232$, $0.9969$ and $2.92\%$. The spreads describe three seeds on the standard
protocol and three initializations on the out-of-distribution protocol.

The $R^2$ convention matters here. The reference weights the three coefficients equally, whereas a
variance-weighted $R^2$ is the default of many libraries. We report the uniform average, computed
from per-coefficient target scales that are recovered from the per-coefficient RMSE of the heads
and checked against the variance-weighted value. The ensemble without correction is close to the
corrected network, while the raw-state kernel is much less accurate in these configurations.

\paragraph{OCO-2 coefficient and radiance objectives.}
Table~\ref{tab:v2-oco2} covers three bands and seeds 0--2. The feature kernel trained under the
flat coefficient loss has coefficient errors of about $4.1\%$, $16.3\%$ and $8.1\%$ on O$_2$, weak
CO$_2$ and strong CO$_2$, and radiance errors of $0.076\%$, $0.114\%$ and $0.108\%$. A
radiance-weighted feature construction gives about $0.031\%$, $0.048\%$ and $0.064\%$ in radiance,
at the price of larger coefficient errors. The reconstruction map and the evaluation metric decide
which fitted representation is preferred.

The radiance stack minimizes a row-relative squared radiance loss, a convex quadratic because the
reconstruction is affine, and its solver reaches KKT gaps between $4.17\times10^{-6}$ and
$1.52\times10^{-5}$. Its mean radiance errors are $0.0264\%$, $0.0391\%$ and $0.0511\%$, against
$0.0291\%$, $0.0404\%$ and $0.0536\%$ for the squared-loss stack on reduced coordinates, while its
coefficient errors rise to $18.59\%$ and $37.43\%$ on O$_2$ and weak CO$_2$ and to $8.87\%$ on
strong CO$_2$. All rules use 2{,}000 validation and 2{,}000 test rows.
